Compartmental modeling is a well-established practice in epidemiological studies. Since the constructs are widely familiar, authors often leave parts of the proposed model implicit when they match usual conventions. As a result, the model described in a paper is not always the same as what is actually validated in the corresponding study, which may inhibit reusability, reproducibility, verifiability and the extensibility of models across diseases and variants. we propose a model-driven completeness auditing framework in which large language models (LLMs) act as extraction instruments for simulated paper-to-model reproduction. We analyzed 10 peer-reviewed modeling papers for 10 infectious diseases (cholera, COVID-19, dengue, Ebola, influenza, HIV, malaria, measles, tuberculosis, and Zika). Across all diseases, using the best-scoring reproduction per paper after comparing three LLM backends (OpenAI, Gemini, and Claude), the mean completeness is $74.0\%$. We discuss what this implies for how epidemiological modeling papers specify and validate models, and for using LLM-assisted tools as systematic instruments for studying completeness rather than as black-box generators.

